\subsection*{Propositional Compactness}

% BEGIN GENERATED ARTIFACT: thm:propositional-compactness
\begin{theorembox}{Theorem (Propositional Compactness)}
\begin{theorem}[Propositional Compactness]
\label{thm:propositional-compactness}
Let \(\Gamma\subseteq\WFF\). If every finite subset of \(\Gamma\) is
satisfiable, then \(\Gamma\) is satisfiable.

\medskip
\noindent\hyperref[prf:propositional-compactness]{Go to proof.}
\end{theorem}
\end{theorembox}

\begin{remark*}[Standard quantified statement]
\[
\operatorname{FinitelySatisfiableSet}(\Gamma)
\Rightarrow
\operatorname{SatisfiableSet}(\Gamma).
\]
\end{remark*}

\begin{remark*}[Predicate reading]
\[
\operatorname{FinitelySatisfiableSet}(\Gamma)
\Rightarrow
\operatorname{SatisfiableSet}(\Gamma)
\]
means that passing every finite satisfiability test is enough to produce a
truth assignment for all of \(\Gamma\).
\end{remark*}

\begin{remark*}[Negated quantified statement]
\[
\operatorname{FinitelySatisfiableSet}(\Gamma)\land
\neg\operatorname{SatisfiableSet}(\Gamma)
\]
is the direct negation of the compactness implication.
\end{remark*}

\begin{remark*}[Negation predicate reading]
\[
\operatorname{FinitelySatisfiableSet}(\Gamma)\land
\neg\operatorname{SatisfiableSet}(\Gamma)
\]
would mean that all finite tests pass but the full set has no satisfying truth
assignment.
\end{remark*}

\begin{remark*}[Failure modes]
\begin{description}
  \item[Exposition.] \[
\neg\operatorname{SatisfiableSet}(\Gamma)
\Rightarrow
\exists\Delta\,\Bigl(\Delta\subseteq\Gamma\land \Delta\text{ is finite}\land
\neg\operatorname{SatisfiableSet}(\Delta)\Bigr).
\]
Compactness says that unsatisfiability cannot first appear only at the infinite
level.
\end{description}
\end{remark*}

\begin{remark*}[Failure modes]
\begin{description}
  \item[Exposition.] \[
\operatorname{UnsatisfiableSet}(\Gamma)
\Rightarrow
\exists\Delta\,\operatorname{FiniteUnsatisfiabilityWitness}(\Delta,\Gamma).
\]
\end{description}
\end{remark*}

\begin{remark*}[Contrapositive quantified statement]
\[
\operatorname{UnsatisfiableSet}(\Gamma)
\Rightarrow
\exists\Delta\,\operatorname{FiniteUnsatisfiabilityWitness}(\Delta,\Gamma).
\]
\end{remark*}

\begin{remark*}[Contrapositive predicate reading]
The contrapositive says that if \(\Gamma\) is unsatisfiable, then the failure is
already witnessed by some finite subset of \(\Gamma\).
\end{remark*}

\begin{remark*}[Interpretation]
Propositional compactness turns infinitely many finite checks into one global
truth assignment. Equivalently, every propositional inconsistency has a finite
witness.
\end{remark*}

\begin{dependencies}
\begin{itemize}
  \item \hyperref[def:finitely-satisfiable-set-of-propositional-formulas]{Finitely Satisfiable Set of Propositional Formulas}
  \item \hyperref[def:satisfiable-set-of-propositional-formulas]{Satisfiable Set of Propositional Formulas}
  \item \hyperref[def:finite-unsatisfiability-witness-propositional-logic]{Finite Unsatisfiability Witness}
\end{itemize}
\end{dependencies}
% END GENERATED ARTIFACT: thm:propositional-compactness
